
\paragraph{Fibonacci 素数场上的“标量化”与 Pisano 周期封口}\label{par:gut-fibprime-pgl2-pisano}
在域相窗口（推论 \ref{cor:field-phase-fib-prime}）中，当模数为 Fibonacci 素数 \(p=F_n\) 时，
Fibonacci 基矩阵
\[
F:=\begin{pmatrix}1&1\\[2pt]1&0\end{pmatrix}
\]
在 \(\mathbb{F}_p\) 上呈现出一种可审计的刚性：其 \(n\) 次幂在 \(PGL_2(\mathbb{F}_p)\) 中回归为单位元，
而在 \(GL_2(\mathbb{F}_p)\) 中仅留下一个由 \(F_{n+1}\) 控制的四次标量因子。
这使得指数 \(n\)、Pisano 周期 \(\pi(p)\) 与若干 Legendre 符号出现严格联动，并与判别式 \(-15\) 的相位锁定二次式
\(2z^2+z+2\) 的跨素数分裂指纹（命题 \ref{prop:pom-phase-lock-discriminant-minus15-splitting}）对齐。

\begin{proposition}[Fibonacci 素数的“标量化”现象：指数 \(n\) 以 \(PGL_2(\mathbb{F}_p)\) 阶的形式回归]\label{prop:gut-fibprime-pgl2-order-n}
设 \(p=F_n\) 为 Fibonacci 素数，并假设 \(n\) 为奇素数（等价地 \(p\ge 5\) 且 \(n\) 为素数指标）。
则在 \(\mathbb{F}_p\) 上有严格同余
\[
F^n\equiv F_{n+1}\,I\pmod p,
\]
从而 \(F\) 在 \(PGL_2(\mathbb{F}_p)\) 的像具有精确阶 \(n\)。
\end{proposition}
\begin{proof}
由标准恒等式
\[
F^n=
\begin{pmatrix}
F_{n+1} & F_n\\
F_n & F_{n-1}
\end{pmatrix},
\]
模 \(p=F_n\) 约化得
\[
F^n\equiv
\mathrm{diag}(F_{n+1},F_{n-1})
\pmod p.
\]
又因 \(F_{n+1}-F_{n-1}=F_n\equiv 0\pmod p\)，故 \(\mathrm{diag}(F_{n+1},F_{n-1})\equiv F_{n+1}I\)。
因此 \(F^n\) 在 \(PGL_2(\mathbb{F}_p)\) 中为单位元，故其像的阶整除 \(n\)。
\par
若该阶小于 \(n\)，由于 \(n\) 为素数只能为 \(1\)，即 \(F\) 在 \(PGL_2(\mathbb{F}_p)\) 中为单位元，等价于 \(F\) 在 \(GL_2(\mathbb{F}_p)\) 中为标量矩阵。
但当 \(p\ge 5\) 时 \(F\) 的非对角元为 \(1\not\equiv 0\pmod p\)，故 \(F\) 不可能为标量矩阵，矛盾。
故其阶必为 \(n\)。证毕。
\leanverified{paper\_gut\_fibprime\_pgl2\_order\_n}
\end{proof}

\begin{corollary}[显式四次剩余：\(F_{n+1}\) 是 \(\mathbb{F}_p\) 中的 \(\sqrt{-1}\)]\label{cor:gut-fibprime-explicit-i}
\leanverified{paper\_gut\_fibprime\_explicit\_i}
在命题 \ref{prop:gut-fibprime-pgl2-order-n} 的设定下，有
\[
F_{n+1}^2\equiv -1\pmod p.
\]
因此 \(\left(\frac{-1}{p}\right)=1\) 且 \(p\equiv 1\pmod 4\)。
更强地，\(\mathbb{F}_p\) 中一个显式的虚单位可取
\[
i\equiv F_{n+1}\pmod p,
\qquad i^2\equiv -1\pmod p.
\]
并且由于 Lucas 数固定为迹变量 \(L_k:=\mathrm{Tr}(F^k)\)，有恒等式
\(
L_n=F_{n-1}+F_{n+1}
\)
从而亦可写为
\[
i\equiv \tfrac12\,L_n\pmod p.
\]
\end{corollary}
\begin{proof}
Cassini 恒等式给出
\[
F_{n+1}F_{n-1}-F_n^2=(-1)^n=-1
\qquad(n\ \text{为奇数}).
\]
模 \(p=F_n\) 化简得 \(F_{n+1}F_{n-1}\equiv -1\pmod p\)。
再用 \(F_{n+1}\equiv F_{n-1}\ (\bmod\ F_n)\) 得 \(F_{n+1}^2\equiv -1\pmod p\)，从而得到显式虚单位。
最后一式由 \(L_n=F_{n-1}+F_{n+1}\) 与 \(F_{n-1}\equiv F_{n+1}\ (\bmod\ p)\) 立得，且 \(p\ge 5\) 保证 \(2\) 在 \(\mathbb{F}_p\) 中可逆。证毕。
\leanverified{paper\_fibprime\_explicit\_sqrt\_neg\_one}
\end{proof}

\begin{corollary}[Pisano 周期的精确封口：\texorpdfstring{$\pi(p)=4n$}{pi(p)=4n}]\label{cor:gut-fibprime-pisano-4n}
设 \(p=F_n\) 为 Fibonacci 素数且 \(n\) 为奇素数。
记 \(\pi(p)\) 为模 \(p\) 的 Pisano 周期。
则有
\[
\pi(p)=4n.
\]
\leanverified{pisano\_period\_17}
\leanverified{pisano\_period\_19}
\leanverified{pisano\_period\_23}
\leanverified{pisano\_wall\_divisibility}
\leanverified{paper\_pisano\_extended}
\leanverified{pisano\_period\_29}
\leanverified{paper\_pisano\_period\_29}
\leanverified{pisano\_period\_31}
\leanverified{pisano\_period\_47}
\leanverified{paper\_pisano\_period\_31\_47}
\leanverified{pisano\_period\_41}
\leanverified{pisano\_period\_59}
\leanverified{paper\_pisano\_period\_41\_59}
\leanverified{pisano\_period\_61}
\leanverified{pisano\_period\_71}
\leanverified{paper\_pisano\_period\_61\_71}
\leanverified{pisano\_period\_89}
\leanverified{paper\_pisano\_period\_89}
\leanverified{pisano\_period\_101}
\leanverified{paper\_pisano\_period\_101}
\leanverified{pisano\_period\_107}
\leanverified{pisano\_period\_113}
\leanverified{paper\_pisano\_period\_107\_113}
\leanverified{paper\_fibprime\_pisano\_rigidity\_package}
\leanverified{paper\_gut\_fibprime\_pisano\_4n}
\leanverified{paper\_pisano\_period\_seeds}
\leanverified{paper\_wall\_seeds}
\end{corollary}
\begin{proof}
经典表示表明 Pisano 周期等于 \(F\) 在 \(GL_2(\mathbb{F}_p)\) 中的阶：
若 \((F_t,F_{t+1})\equiv(0,1)\ (\bmod\ p)\)，则由矩阵幂恒等式
\(
F^t=\bigl(\begin{smallmatrix}F_{t+1}&F_t\\F_t&F_{t-1}\end{smallmatrix}\bigr)
\)
得 \(F^t\equiv I\)；
反向地 \(F^t\equiv I\) 立即推出 \((F_t,F_{t+1})\equiv(0,1)\)。
\par
由命题 \ref{prop:gut-fibprime-pgl2-order-n} 有 \(F^n\equiv F_{n+1}I\)；
由推论 \ref{cor:gut-fibprime-explicit-i} 有 \(F_{n+1}^2\equiv -1\)，故
\[
F^{2n}\equiv -I,\qquad F^{4n}\equiv I.
\]
由于 \(p\ge 5\)，\(-I\not\equiv I\)，从而 \(F^{2n}\not\equiv I\)；因此 \(\mathrm{ord}_{GL_2(\mathbb{F}_p)}(F)\) 不可能整除 \(2n\)，而又整除 \(4n\)，于是其阶必为 \(4n\)。
故 \(\pi(p)=4n\)。证毕。
\leanverified{paper\_fibprime\_pisano\_4n}
\end{proof}

\begin{proposition}[指数—素数值的强同余刚性：\texorpdfstring{$p\equiv (\frac{5}{p})\ (\bmod\ n)$}{p ≡ (5/p) mod n}]\label{prop:gut-fibprime-congruence-p-mod-n}
设 \(p=F_n\) 为 Fibonacci 素数，且 \(n\) 为奇素数（\(n>5\)）。
则必有
\[
p \equiv \left(\frac{5}{p}\right)\pmod n,
\]
从而 \(p\equiv 1\ (\bmod\ n)\) 或 \(p\equiv -1\ (\bmod\ n)\)，并且符号由 Legendre 符号 \(\left(\frac{5}{p}\right)\in\{\pm 1\}\) 唯一决定。
\end{proposition}
\begin{proof}
令 \(r(p):=\min\{k\ge 1:\ p\mid F_k\}\) 为 Fibonacci 序列在模 \(p\) 下的出现阶（rank of apparition）。
对任意奇素数 \(p\neq 5\)，经典 Lucas/Fibonacci 整除律给出
\[
r(p)\ \mid\ \Bigl(p-\left(\frac{5}{p}\right)\Bigr).
\]
在本命题中 \(p=F_n\) 且 \(n\) 为素数指标；由 Fibonacci 严格递增可知对一切 \(k<n\) 有 \(F_k<p\)，故 \(p\nmid F_k\)（\(k<n\)），从而 \(r(p)=n\)。
代回即得 \(n\mid\bigl(p-(\frac{5}{p})\bigr)\)，等价于所述同余。证毕。
\leanverified{paper\_gut\_fibprime\_congruence\_p\_mod\_n}
\end{proof}

\begin{corollary}[符号的显式化：由 \texorpdfstring{$n\ (\mathrm{mod}\ 20)$}{n (mod 20)} 决定 \texorpdfstring{$p\equiv\pm 1\ (\mathrm{mod}\ n)$}{p ≡ ±1 (mod n)}]\label{cor:gut-fibprime-sign-by-n-mod20}
设 \(p=F_n\) 为 Fibonacci 素数，且 \(n\) 为奇素数（\(n>5\)）。
则
\[
n\equiv 1,9,11,19\pmod{20}\ \Longrightarrow\ \left(\frac{5}{p}\right)=1\ \Longrightarrow\ p\equiv 1\pmod n,
\]
\[
n\equiv 3,7,13,17\pmod{20}\ \Longrightarrow\ \left(\frac{5}{p}\right)=-1\ \Longrightarrow\ p\equiv -1\pmod n.
\]
\end{corollary}
\begin{proof}
\(\left(\frac{5}{p}\right)=1\) 当且仅当 \(p\equiv \pm 1\ (\bmod\ 5)\)，而 \(\left(\frac{5}{p}\right)=-1\) 当且仅当 \(p\equiv \pm 2\ (\bmod\ 5)\)。
另一方面，\(p=F_n\ (\bmod\ 5)\) 仅依赖于 \(n\ (\bmod\ 20)\)（因为 \(\bmod 5\) 的 Pisano 周期为 \(20\)）。
对奇素数 \(n>5\) 的允许余类恰分裂为上述两组，从而确定 \(\left(\frac{5}{p}\right)\) 的符号；
再由命题 \ref{prop:gut-fibprime-congruence-p-mod-n} 得到 \(p\ (\bmod\ n)\) 的符号。证毕。
\leanverified{paper\_fibprime\_sign\_by\_n\_mod20}
\leanverified{paper\_gut\_fibprime\_sign\_by\_n\_mod20}
\end{proof}

\begin{corollary}[与 \texorpdfstring{$A_2$}{A2} 判别式 \texorpdfstring{$-15$}{-15} 的结构对接：Fibonacci 素数场中的相位锁定根判据]\label{cor:gut-fibprime-disc-minus15-criterion}
\leanverified{paper\_gut\_fibprime\_disc\_minus15\_criterion}
设 \(p=F_n\) 为 Fibonacci 素数且 \(n\) 为奇素数。
考虑二次锁定多项式 \(2z^2+z+2\)（判别式 \(\Delta=-15\)）。
则该多项式在 \(\mathbb{F}_p\) 上可分解当且仅当
\[
\left(\frac{15}{p}\right)=1.
\]
在可分解时，其两根可写为
\[
z=\frac{-1\pm i\sqrt{15}}{4}\in\mathbb{F}_p,
\]
其中 \(i\in\mathbb{F}_p\) 可取为推论 \ref{cor:gut-fibprime-explicit-i} 的显式值 \(i\equiv F_{n+1}\ (\bmod\ p)\)。
该判据与命题 \ref{prop:pom-phase-lock-discriminant-minus15-splitting} 的一般分裂选择律一致，并在 Fibonacci 素数场中被推论 \ref{cor:gut-fibprime-explicit-i} 的 \(\left(\frac{-1}{p}\right)=1\) 进一步压缩为单一 Legendre 条件。
\end{corollary}
\begin{proof}
在任意奇素数域 \(\mathbb{F}_p\) 中，二次式可分解当且仅当判别式为二次剩余。
此处 \(\Delta=-15\)，故分解性等价于 \(\left(\frac{-15}{p}\right)=1\)。
由推论 \ref{cor:gut-fibprime-explicit-i} 得 \(\left(\frac{-1}{p}\right)=1\)，从而
\(
\left(\frac{-15}{p}\right)=\left(\frac{-1}{p}\right)\left(\frac{15}{p}\right)=\left(\frac{15}{p}\right)
\)，即得到所述等价。
根式表达由二次公式给出，并以 \(i^2=-1\) 实现 \(\sqrt{-15}=i\sqrt{15}\)；而 \(i\equiv F_{n+1}\) 已由推论 \ref{cor:gut-fibprime-explicit-i} 给出。证毕。
\leanverified{paper\_fibprime\_disc\_minus15\_criterion}
\end{proof}

\endinput
